For parameters $q,t\in(0,1)$ and a partition $\lambda$ with at most $n$ parts, $P_\lambda(x;q,t)$ denotes the Macdonald polynomial in the monic normalization of Macdonald's book, as used in \citet{McSwiggenSahi2026}; at $q=t$ it specializes to the Schur polynomial $s_\lambda(x)$, which is what the refutation below exploits.  The statement concerns two convexity properties of the map $\lambda\mapsto\Omega_\lambda(x;q,t)$ on partitions: \emph{log-convexity}, and \emph{Schur-convexity}, meaning monotonicity with respect to the dominance (majorization) order, so that $\lambda\succeq\mu$ implies $\Omega_\lambda\ge\Omega_\mu$.  The normalization $\Omega_\lambda$ and the lattice $\mathcal L_n$ of evaluation points are defined inside the quoted statement.
